\begin{lemma}[Existence of the complex $L^1$ extension satisfying the $8/\pi$ weak bound]
There is a $\mathbb{C}$-linear map
\[
  T_{\mathbb{C}} : L^1(\gamma_n;\mathbb{C}) \longrightarrow
    \mathcal{M}(\gamma_n;\mathbb{C}^n),
\]
where $\mathcal{M}(\gamma_n;\mathbb{C}^n)$ denotes the space of
$\gamma_n$-almost-everywhere equivalence classes of measurable
$\mathbb{C}^n$-valued functions on $\mathbb{R}^n$, such that
\begin{enumerate}
\item $T_{\mathbb{C}}$ agrees with $R_\gamma^{\mathbb{C}}$ on
  $L^2(\gamma_n;\mathbb{C})$: for every $v \in L^2(\gamma_n;\mathbb{C})$ one has
  $T_{\mathbb{C}}(\iota_{\mathbb{C}} v) = R_\gamma^{\mathbb{C}} v$, with
  $\iota_{\mathbb{C}}$ the inclusion of \noderef{ComplexToL1} and
  $R_\gamma^{\mathbb{C}}$ as in \noderef{ComplexRieszTransform}; and
\item $T_{\mathbb{C}}$ satisfies the weak bound
\[
  t\,\gamma_n\bigl(\{x : |T_{\mathbb{C}}f(x)|_{\mathbb{C}^n} > t\}\bigr)
    \leq \frac{8}{\pi}\,\|f\|_{L^1(\gamma_n;\mathbb{C})}
  \qquad\text{for every } f \in L^1(\gamma_n;\mathbb{C}) \text{ and } t > 0 .
\]
\end{enumerate}

This is the existence half of the reference paper's last sentence, that
approximation by complex $L^2$ functions gives a continuous complex-linear
extension to $L^1$ --- ``continuous'' being clause (2), which is boundedness
from $L^1(\gamma_n;\mathbb{C})$ into weak $L^1$ with quasi-norm at most $8/\pi$.
The scalars in the linearity are the complex ones, which is what makes the map
an extension of the complex transform rather than merely of its real and
imaginary parts.
\end{lemma}
\begin{proof}
\emph{Step 1: the definition of $T_{\mathbb{C}}$, and its characterizing
property.}
Let $f \in L^1(\gamma_n;\mathbb{C})$.  Apply \noderef{ComplexRieszL1Limit} to
$f$ and let $g_f \in \mathcal{M}(\gamma_n;\mathbb{C}^n)$ and
$(v^f_k) \subseteq L^2(\gamma_n;\mathbb{C})$ be the element and the sequence it
supplies; define $T_{\mathbb{C}}f = g_f$.  Clause (4) of that node then reads:
\begin{equation*}
  \text{for every } (w_k) \subseteq L^2(\gamma_n;\mathbb{C}) \text{ with }
  \iota_{\mathbb{C}} w_k \to f \text{ in } L^1,
  \quad R_\gamma^{\mathbb{C}} w_k \to T_{\mathbb{C}}f \text{ in measure.}
  \tag{$\star$}
\end{equation*}
Property $(\star)$ determines $T_{\mathbb{C}}f$ as an element of
$\mathcal{M}(\gamma_n;\mathbb{C}^n)$: clause (2) of the same node supplies at
least one such sequence, namely $(v^f_k)$, and a limit in measure is unique up to
a $\gamma_n$-null set, because if $A_k \to A$ and $A_k \to A'$ in measure then for
every $\varepsilon > 0$,
$\gamma_n(\{|A - A'| \geq 2\varepsilon\})
 \leq \gamma_n(\{|A - A_k| \geq \varepsilon\})
    + \gamma_n(\{|A_k - A'| \geq \varepsilon\}) \to 0$, so that measure is $0$
for every $\varepsilon$, and $\{A \neq A'\}$ is the countable union of the sets
$\{|A-A'| \geq 2/j\}$.  Everything below uses only $(\star)$ together with
clauses (1)--(3) at the particular sequence $(v^f_k)$.

\emph{Step 2: convergence in measure respects the vector operations.}
If $A_k \to A$ and $B_k \to B$ in measure and $c \in \mathbb{C}$, then
$A_k + B_k \to A + B$ and $c\,A_k \to c\,A$ in measure.  For the sum: if
$|A_k + B_k - (A+B)| \geq \varepsilon$ then $|A_k - A| \geq \varepsilon/2$ or
$|B_k - B| \geq \varepsilon/2$, so by monotonicity and subadditivity of
$\gamma_n$ (shared context) the measure of the first set is at most the sum of
the measures of the other two, both of which tend to $0$.  For the scalar
multiple: if $c = 0$ both sides are $0$; otherwise
$\{|c A_k - cA| \geq \varepsilon\} = \{|A_k - A| \geq \varepsilon/|c|\}$, whose
measure tends to $0$.

\emph{Step 3: $\mathbb{C}$-linearity.}
Let $f, h \in L^1(\gamma_n;\mathbb{C})$ and $c \in \mathbb{C}$, and let
$(v^f_k), (v^h_k)$ be the sequences of Step 1 for $f$ and for $h$.  Since
$\iota_{\mathbb{C}}$ is $\mathbb{C}$-linear (\noderef{ComplexToL1}),
$\iota_{\mathbb{C}}(v^f_k + v^h_k)
 = \iota_{\mathbb{C}} v^f_k + \iota_{\mathbb{C}} v^h_k$, and by the triangle
inequality in $L^1$,
\[
  \bigl\|\iota_{\mathbb{C}}(v^f_k + v^h_k) - (f+h)\bigr\|_{L^1}
    \leq \|\iota_{\mathbb{C}} v^f_k - f\|_{L^1}
      + \|\iota_{\mathbb{C}} v^h_k - h\|_{L^1} \longrightarrow 0
\]
by clause (2) of \noderef{ComplexRieszL1Limit} at $f$ and at $h$; so
$\iota_{\mathbb{C}}(v^f_k + v^h_k) \to f + h$ in $L^1$, and
$(\star)$ at $f+h$ gives $R_\gamma^{\mathbb{C}}(v^f_k + v^h_k)
 \to T_{\mathbb{C}}(f+h)$ in measure.  On the other hand clause (1) of
\noderef{ComplexRieszTransformLinear} gives
$R_\gamma^{\mathbb{C}}(v^f_k + v^h_k)
 = R_\gamma^{\mathbb{C}} v^f_k + R_\gamma^{\mathbb{C}} v^h_k$, and by $(\star)$ at
$f$ and at $h$ the two summands converge in measure to $T_{\mathbb{C}}f$ and
$T_{\mathbb{C}}h$, so by Step 2 their sum converges in measure to
$T_{\mathbb{C}}f + T_{\mathbb{C}}h$.  Uniqueness of limits in measure (Step 1)
gives $T_{\mathbb{C}}(f+h) = T_{\mathbb{C}}f + T_{\mathbb{C}}h$.

Likewise $\iota_{\mathbb{C}}(c\,v^f_k) = c\,\iota_{\mathbb{C}} v^f_k$ by
$\mathbb{C}$-linearity, and
$\|c\,\iota_{\mathbb{C}} v^f_k - c f\|_{L^1}
 = |c|\,\|\iota_{\mathbb{C}} v^f_k - f\|_{L^1} \to 0$ by absolute homogeneity
of the $L^1$ norm, so $\iota_{\mathbb{C}}(c\,v^f_k) \to c f$ in $L^1$ and $(\star)$ at $cf$ gives $R_\gamma^{\mathbb{C}}(c\,v^f_k)
 \to T_{\mathbb{C}}(cf)$ in measure; clause (2) of
\noderef{ComplexRieszTransformLinear} gives
$R_\gamma^{\mathbb{C}}(c\,v^f_k) = c\,R_\gamma^{\mathbb{C}} v^f_k$, which
converges in measure to $c\,T_{\mathbb{C}}f$ by $(\star)$ at $f$ and Step 2.
Hence $T_{\mathbb{C}}(cf) = c\,T_{\mathbb{C}}f$.  So $T_{\mathbb{C}}$ is a
$\mathbb{C}$-linear map $L^1(\gamma_n;\mathbb{C}) \to
\mathcal{M}(\gamma_n;\mathbb{C}^n)$.

\emph{Step 4: clause (1), agreement on $L^2$.}
Let $v \in L^2(\gamma_n;\mathbb{C})$ and take the constant sequence
$w_k = v$.  Then $\iota_{\mathbb{C}} w_k = \iota_{\mathbb{C}} v$ for every $k$,
so trivially $\iota_{\mathbb{C}} w_k \to \iota_{\mathbb{C}} v$ in $L^1$, and
$(\star)$ at $f = \iota_{\mathbb{C}} v$ says that the constant sequence
$R_\gamma^{\mathbb{C}} v$ converges in measure to
$T_{\mathbb{C}}(\iota_{\mathbb{C}} v)$.  A constant sequence converges in measure
to its own value, so by uniqueness of limits in measure (Step 1)
$T_{\mathbb{C}}(\iota_{\mathbb{C}} v) = R_\gamma^{\mathbb{C}} v$.

\emph{Step 5: clause (2), the weak bound.}
Fix $f \in L^1(\gamma_n;\mathbb{C})$ and $t > 0$, and let $(v^f_k)$ be the
sequence of Step 1, so that by clause (3) of \noderef{ComplexRieszL1Limit} we
have $R_\gamma^{\mathbb{C}} v^f_k \to g_f = T_{\mathbb{C}}f$ pointwise
$\gamma_n$-a.e.  At a point $x$ with $|T_{\mathbb{C}}f(x)| > t$ the margin
$|T_{\mathbb{C}}f(x)| - t$ is strictly positive and
$|R_\gamma^{\mathbb{C}} v^f_k(x)| \to |T_{\mathbb{C}}f(x)|$ by continuity of the
norm, so $|R_\gamma^{\mathbb{C}} v^f_k(x)| > t$ for all large $k$.  Hence
$\gamma_n$-a.e.
\[
  \mathbf{1}_{\{|T_{\mathbb{C}}f| > t\}}
    \leq \liminf_{k\to\infty}\mathbf{1}_{\{|R_\gamma^{\mathbb{C}} v^f_k| > t\}} ,
\]
and Fatou's lemma (shared context), applied to these nonnegative measurable
functions, gives
\[
  \gamma_n\bigl(\{|T_{\mathbb{C}}f| > t\}\bigr)
    \leq \liminf_{k}\,
      \gamma_n\bigl(\{|R_\gamma^{\mathbb{C}} v^f_k| > t\}\bigr) .
\]
By \noderef{ComplexMainBound} applied to $v^f_k$ at the level $t > 0$, and then
by $\|\iota_{\mathbb{C}} v^f_k\|_{L^1} = \int |v^f_k| \, d\gamma_n$
(\noderef{ComplexToL1}) followed by clause (1) of
\noderef{ComplexRieszL1Limit},
\[
  t\,\gamma_n\bigl(\{|R_\gamma^{\mathbb{C}} v^f_k| > t\}\bigr)
    \leq \frac{8}{\pi}\int |v^f_k| \, d\gamma_n
    = \frac{8}{\pi}\,\|\iota_{\mathbb{C}} v^f_k\|_{L^1}
    \leq \frac{8}{\pi}\,\|f\|_{L^1} ,
\]
so every term of the $\liminf$ is at most $\tfrac{8}{\pi t}\|f\|_{L^1}$ and
therefore
$\gamma_n(\{|T_{\mathbb{C}}f| > t\}) \leq \tfrac{8}{\pi t}\|f\|_{L^1}$.
Multiplying by $t > 0$ gives clause (2).
\end{proof}
